\VerdierPhysicalPageAnchor{0048}%
\VerdierSemanticAnchor{ch1-1-4-4}
\noindent\textbf{Definition 1.4.4.}\quad
Let $(\mathcal C,G,T,c)$ and $(\mathcal C',G,T',c')$ be two
$G$-categories, with $\mathcal C'$ an additive category. Let
$\alpha\in Z^2(G,\mathbb Z/2\mathbb Z)$ be a normalized
inhomogeneous cocycle. A \emph{$G$-functor twisted by the cocycle
$\alpha$} is an object $(F,p)$ consisting of a functor

\VerdierSemanticAnchor{ch1-1-4-4-functor}
\[
  F\colon\mathcal C\longrightarrow\mathcal C'
\]
and functor isomorphisms

\VerdierSemanticAnchor{ch1-1-4-4-twisting-data}
\[
  p(\sigma)\colon FT(\sigma)\longrightarrow T'(\sigma)F,
  \qquad \sigma\in G,
\]
such that, for every pair $\sigma_1,\sigma_2\in G$, the diagram
(1.4.3.11) is $\alpha(\sigma_1,\sigma_2)$-commutative.

\VerdierSemanticAnchor{ch1-1-4-4-correspondence}
The choice of a normalized section $s$ of the extension (1.4.3.4)
therefore establishes a bijective correspondence between the
$H$-functors
\[
  (F,m)\colon(\mathcal C,H,T_f,c_f)
  \longrightarrow(\mathcal C',H,T'_f,c'_f)
\]
that induce the canonical operation of the group of signs on $F$, and
the $G$-functors twisted by the cocycle $\alpha_s$ defined by the
section $s$.

\VerdierSemanticAnchor{ch1-1-4-5}
\noindent\textbf{1.4.5.}\quad
We use the notation of (1.4.3). Let

\VerdierSemanticAnchor{ch1-1-4-5-eq1}
\begin{equation}
  \begin{gathered}
    (F_1,m_1),(F_2,m_2)\colon
    (\mathcal C,H,T_f,c_f)\\[-0.2em]
    \longrightarrow(\mathcal C',H,T'_f,c'_f)
  \end{gathered}
  \tag{1.4.5.1}
\end{equation}
be two $H$-functors inducing the canonical operation of the group of
signs, and let

\VerdierSemanticAnchor{ch1-1-4-5-eq2}
\begin{equation}
  u\colon(F_1,m_1)\longrightarrow(F_2,m_2)
  \tag{1.4.5.2}
\end{equation}
be a morphism of functors compatible with the operations of $H$.

\VerdierSemanticAnchor{ch1-1-4-5-sections}
Let $s_1$ and $s_2$ be two normalized sections of the map $f$ in the
exact sequence (1.4.3.4) ($s_1(e)=s_2(e)=e$). Let
$\alpha_{s_1},\alpha_{s_2}\in Z^2(G,\mathbb Z/2\mathbb Z)$ be the
corresponding normalized cocycles. Let
$\beta\colon G\to\mathbb Z/2\mathbb Z$ be the map

\VerdierSemanticAnchor{ch1-1-4-5-eq3}
\begin{equation}
  \beta(\sigma)=s_2(\sigma)s_1(\sigma)^{-1},
  \qquad \sigma\in G.
  \tag{1.4.5.3}
\end{equation}

\VerdierSemanticAnchor{ch1-1-4-5-induced}
By (1.4.3), the functors $(F_1,m_1)$ and $(F_2,m_2)$ define
$G$-functors $(F_1,p_1)$ and $(F_2,p_2)$ twisted respectively by the
cocycles $\alpha_{s_1}$ and $\alpha_{s_2}$. One verifies immediately
that the morphism $u\colon F_1\to F_2$ between the underlying functors
has the following property:

\VerdierSemanticAnchor{ch1-1-4-5-property}
